% ============================================================
% CuFun: Theoretical Framework
% ============================================================

\section{Theoretical Framework}
\label{sec:theory}

\subsection{Fundamental Assumptions}

\begin{assumption}[Temporal Point Process Foundation]
\label{assumption:foundation}
The observed event sequence $\{t_1, t_2, \ldots, t_n\}$ is generated by an underlying
temporal point process. The event history $\mathcal{H}_{i-1}$ can be adequately
summarized by a finite-dimensional embedding vector $\boldsymbol{h}_{i-1} \in \mathcal{H}
\subset \mathbb{R}^d$ (compact) such that inter-event times depend only on this embedding.
\end{assumption}

\begin{assumption}[CDF Regularity]
\label{assumption:regularity}
The true conditional CDF $F_0(\tau|\boldsymbol{h})$ is continuous, strictly increasing
for $\tau > 0$, and satisfies $F_0(0|\boldsymbol{h}) = 0$,
$\lim_{\tau \to \infty} F_0(\tau|\boldsymbol{h}) = 1$ for all $\boldsymbol{h} \in \mathcal{H}$.
\end{assumption}

\begin{assumption}[Joint H\"{o}lder Smoothness]
\label{assumption:smoothness}
$F_0$ belongs to the $\alpha$-H\"{o}lder class \emph{jointly} in $(\tau, \boldsymbol{h})$:
there exists $C > 0$ such that for all $\tau_1, \tau_2 \geq 0$ and
$\boldsymbol{h}_1, \boldsymbol{h}_2 \in \mathcal{H}$,
\[
|F_0(\tau_1|\boldsymbol{h}_1) - F_0(\tau_2|\boldsymbol{h}_2)|
\;\leq\; C\bigl(|\tau_1-\tau_2|^\alpha + \|\boldsymbol{h}_1 - \boldsymbol{h}_2\|^\alpha\bigr).
\]
In particular, $F_0(\cdot\,|\boldsymbol{h})$ is $\alpha$-H\"{o}lder in $\tau$ uniformly
over $\boldsymbol{h}$, and $F_0(\tau|\cdot)$ is $\alpha$-H\"{o}lder in $\boldsymbol{h}$
uniformly over $\tau$.
\end{assumption}

\begin{assumption}[Network Architecture]
\label{assumption:architecture}
The MNN $\mathcal{M}(\tau, \boldsymbol{h}; \theta)$ has $\geq 2$ hidden layers,
activations $\phi \in C^1$ that are non-decreasing and satisfy $\phi' > 0$
on a dense subset, and all weight matrices $W^{(\ell)}_\tau$ on computational
paths from $\tau$ are constrained to be entry-wise non-negative.
\end{assumption}

\begin{assumption}[Conditional Independence]
\label{assumption:conditional_independence}
Given $\boldsymbol{h}_{i-1}$, $\tau_i$ is conditionally independent of
$\{\tau_j\}_{j < i-1}$, with $\tau_i \sim F_0(\cdot|\boldsymbol{h}_{i-1})$.
\end{assumption}

% ============================================================
\subsection{MNN Structural Properties}
% ============================================================

\begin{theorem}[MNN Architectural Monotonicity and Boundedness \verified{}]
\label{thm:mnn-mono}
Under Assumption~\ref{assumption:architecture}, the MNN output
$F^*(\tau|\boldsymbol{h}) := \sigma\!\bigl(\mathcal{M}(\tau,\boldsymbol{h};\theta)\bigr)$
is (i) monotone non-decreasing in $\tau$, (ii) bounded in $(0,1)$, and
(iii) has well-defined log-density $\log(dF^*/d\tau) > -\infty$.
\end{theorem}

\begin{proof}
(i)~By the non-negative weight constraint (Assumption~\ref{assumption:architecture}),
$\mathcal{M}(\tau_2)-\mathcal{M}(\tau_1) = \sum_i w_i(\phi(\tau_2)-\phi(\tau_1)) \geq 0$
for $\tau_1 \leq \tau_2$; since $\sigma$ is strictly increasing, $F^*(\tau_1)\leq F^*(\tau_2)$.
(ii)~$\sigma(z)\in(0,1)$ for all $z\in\R$.
(iii)~$dF^*/d\tau = \sigma'(\mathcal{M})\cdot\mathcal{M}'_\tau$;
$\sigma'(z) = \sigma(z)(1-\sigma(z)) > 0$ everywhere, and
$\mathcal{M}'_\tau \geq 0$ by the non-negative weight constraint,
with strict positivity on the dense set $\{\phi'>0\}$ guaranteed by
Assumption~\ref{assumption:architecture}.
\end{proof}

% ============================================================
\subsection{MNN Structural Guarantees}
\label{subsec:mnn}
% ============================================================

\begin{theorem}[MNN Monotonicity and Boundedness]
\label{thm:mnn-mono}
Under Assumption~\ref{assumption:architecture}, the MNN output
$F^*(\tau|\boldsymbol{h}) = \sigma(\mathcal{M}(\tau,\boldsymbol{h};\theta))$
is non-decreasing in $\tau$ and satisfies $F^*(\tau|\boldsymbol{h}) \in (0,1)$
for all $\tau \geq 0$, $\boldsymbol{h} \in \mathcal{H}$, and all $\theta$.
Hence $F^*$ constitutes a valid conditional CDF for every parameter vector.
\end{theorem}

\begin{proof}
\emph{Monotonicity.}
The non-negative weight constraint on the $\tau$-path (Assumption~\ref{assumption:architecture})
ensures $\mathcal{M}(\tau,\boldsymbol{h};\theta)$ is non-decreasing in $\tau$
(a composition of non-decreasing functions with non-negative linear maps is
non-decreasing). Since $\sigma$ is strictly increasing, $F^*$ inherits monotonicity.
\emph{Boundedness.} $\sigma:\mathbb{R}\to(0,1)$ maps any finite pre-activation to $(0,1)$,
so $F^* \in (0,1)$ holds by architecture without any constraint on $\theta$.
\end{proof}

\begin{theorem}[CDF--Intensity Derivative Identity]
\label{thm:cdf-intensity}
Let $F^*(\tau|\boldsymbol{h}) = \sigma(\mathcal{M}(\tau,\boldsymbol{h};\theta))$.
The per-event log-likelihood is given exactly by automatic differentiation:
\[
  \log p^*(\tau|\boldsymbol{h})
  = \log\frac{\partial F^*}{\partial\tau}
  = \log\!\bigl[\sigma'(\mathcal{M})\cdot\mathcal{M}'_\tau\bigr]
  = \log\!\bigl[\sigma(\mathcal{M})(1-\sigma(\mathcal{M}))\bigr]
    + \log\mathcal{M}'_\tau,
\]
where $\sigma'(z) = \sigma(z)(1-\sigma(z)) > 0$ for all $z$, so $\log p^*$ is
always finite and requires no numerical integration.
\end{theorem}

\begin{proof}
Chain rule: $\partial F^*/\partial\tau = \sigma'(\mathcal{M})\cdot\mathcal{M}'_\tau$.
Both factors are positive ($\sigma'(z)>0$ and $\mathcal{M}'_\tau > 0$ by
Assumption~\ref{assumption:architecture}), so the log is finite.
Taking the log gives the identity above. $\square$
\end{proof}

% ============================================================
\subsection{Universal Approximation}
% ============================================================

The non-triviality arises in three places: (i) the domain is \emph{unbounded}
in $\tau$, (ii) the target is a \emph{joint} function of time and history,
and (iii) the approximation must \emph{simultaneously} preserve CDF validity.

\begin{lemma}[Exponential Tail Truncation]
\label{lem:tail}
Under Assumptions~\ref{assumption:regularity} and~\ref{assumption:smoothness},
for every $\epsilon > 0$ there exists $T^* = T^*(\epsilon, \mathcal{H}) < \infty$ such that
\[
\inf_{\boldsymbol{h} \in \mathcal{H}} F_0(T^*|\boldsymbol{h}) \;>\; 1 - \epsilon.
\]
\end{lemma}

\begin{proof}
By Assumption~\ref{assumption:regularity}, $\lim_{\tau\to\infty} F_0(\tau|\boldsymbol{h}) = 1$
for each $\boldsymbol{h}$.
By the joint H\"{o}lder bound (Assumption~\ref{assumption:smoothness}),
for any $\boldsymbol{h}$ with $\|\boldsymbol{h}-\boldsymbol{h}_0\| \leq \delta$:
\[
F_0(\tau|\boldsymbol{h}) \geq F_0(\tau|\boldsymbol{h}_0) - C\delta^\alpha.
\]
Since $\mathcal{H}$ is compact, cover it with finitely many balls $B(\boldsymbol{h}_j, \delta)$.
Choose $\delta = (\epsilon/2C)^{1/\alpha}$ so the perturbation is $< \epsilon/2$.
For each centre $\boldsymbol{h}_j$, choose $T_j$ such that
$F_0(T_j|\boldsymbol{h}_j) > 1 - \epsilon/2$.
Set $T^* = \max_j T_j$; then for all $\boldsymbol{h} \in \mathcal{H}$:
\[
F_0(T^*|\boldsymbol{h}) \geq F_0(T^*|\boldsymbol{h}_j) - C\delta^\alpha
> (1 - \tfrac{\epsilon}{2}) - \tfrac{\epsilon}{2} = 1 - \epsilon. \qed
\]
\end{proof}

\begin{lemma}[Piecewise-Linear Monotone Approximation]
\label{lem:pwl}
For any continuous non-decreasing $g: [0,T] \to [0,1]$ with $g(0) = 0$
and $\alpha$-H\"{o}lder constant $C_g$,
and any $\epsilon > 0$, there exists a non-decreasing piecewise-linear
$\hat{g}: [0,T] \to [0,1]$ with $\hat{g}(0)=0$ and
$\lceil (C_g T^\alpha/\epsilon)^{1/\alpha} \rceil = O(\epsilon^{-1/\alpha})$ breakpoints such that
$\|\hat{g} - g\|_{\infty,[0,T]} < \epsilon$.
Moreover, $\hat{g}$ is realisable as a two-layer ReLU network with non-negative weights.
\end{lemma}

\begin{proof}
Partition $[0,T]$ into $N$ equal subintervals $[s_k, s_{k+1}]$ with $s_k = kT/N$.
Construct the piecewise-linear interpolant $\hat{g}$ through the breakpoints $(s_k, g(s_k))$.
By H\"{o}lder continuity, $\|g - \hat{g}\|_\infty \leq C_g (T/N)^\alpha$.
Setting $N = \lceil (C_g T^\alpha/\epsilon)^{1/\alpha} \rceil$ yields $\|g-\hat{g}\|_\infty < \epsilon$.

The interpolant is realised exactly as:
\[
\hat{g}(\tau) = \sum_{k=0}^{N-1} \underbrace{(g(s_{k+1})-g(s_k))}_{\geq\, 0}
\cdot \operatorname{ReLU}\!\left(\frac{\tau - s_k}{T/N}\right)
\cdot \mathbf{1}[\tau \leq s_{k+1}],
\]
where all coefficients are non-negative (since $g$ is non-decreasing), so the
network has non-negative weights. Since $g(0)=0$, the bias vanishes, and
$\hat{g}(0)=0$.
\end{proof}

\begin{theorem}[Universal Approximation with Quantitative Rates]
\label{thm:universal_approx}
Under Assumptions~\ref{assumption:regularity}--\ref{assumption:architecture}
and~\ref{assumption:smoothness}, for any $\epsilon > 0$, there exists a
network configuration $\theta^*$ with
\[
N^* = O\!\left(\epsilon^{-(1+d/\alpha)}\right) \;\text{ parameters}
\]
such that
\[
\sup_{\tau \geq 0,\, \boldsymbol{h} \in \mathcal{H}}
\bigl|F^*(\tau|\boldsymbol{h}) - F_0(\tau|\boldsymbol{h})\bigr| < \epsilon,
\]
where $F^*(\tau|\boldsymbol{h}) = \sigma\!\bigl(\mathcal{M}(\tau,\boldsymbol{h};\theta^*)\bigr)$
is a valid CDF for all~$\theta^*$.
\end{theorem}

\begin{proof}
We construct a specific network satisfying Assumption~\ref{assumption:architecture}.
The proof proceeds in two stages: \textbf{Stage A} (Steps~1--4) builds a monotone
$[0,1]$-valued function $\hat{\mathcal{M}}$ with
$\|\hat{\mathcal{M}} - F_0\|_\infty < 3\epsilon/8$;
\textbf{Stage B} (Step~5) embeds $\hat{\mathcal{M}}$ into the $\sigma$-MNN
architecture, incurring at most $\epsilon/4 + \epsilon/8$ additional error,
so that the total is $3\epsilon/4 < \epsilon$.

\medskip\noindent\textbf{Preliminary (Gaussian softmax $\alpha$-moment).}
For a $\delta$-net $\{\boldsymbol{h}_j\}_{j=1}^K$ and Gaussian softmax weights
\[
w_j(\boldsymbol{h})
= \frac{\exp(-\|\boldsymbol{h}-\boldsymbol{h}_j\|^2/\delta^2)}
       {\sum_{k} \exp(-\|\boldsymbol{h}-\boldsymbol{h}_k\|^2/\delta^2)},
\quad \sum_j w_j \equiv 1, \quad w_j \geq 0,
\]
define $\Gamma_{d,\alpha}
:= \sup_{\boldsymbol{h}\in\mathcal{H}}\sum_j w_j(\boldsymbol{h})
   (\|\boldsymbol{h}_j-\boldsymbol{h}\|/\delta)^\alpha$.
This constant is finite: centres at normalised distance $r > 1$ receive weight
$w_j \leq e^{1-r^2}$ (since the denominator is $\geq e^{-1}$), and the shell
$[r\delta, (r+1)\delta]$ contains $O(r^{d-1})$ centres, so
$\Gamma_{d,\alpha} \leq 1 + e\sum_{r=1}^\infty C'_d r^{d-1+\alpha} e^{-r^2} < \infty$.
Write $\Gamma$ for brevity.

\medskip\noindent\textsc{Stage A: Direct CDF Approximation.}

\textbf{Step 1 (Tail and boundary reduction).}
Apply Lemma~\ref{lem:tail} with $\epsilon/8$ to choose $T^*$ such that
$\inf_{\boldsymbol{h}} F_0(T^*|\boldsymbol{h}) > 1 - \epsilon/8$.
By continuity and compactness, choose $\tau_{\min} > 0$ with
$\sup_{\boldsymbol{h}} F_0(\tau_{\min}|\boldsymbol{h}) < \epsilon/8$.

\textbf{Step 2 (History discretisation).}
Set $\delta = (\epsilon / (8C\max\{1,\Gamma\}))^{1/\alpha}$.
Cover $\mathcal{H}$ with a $\delta$-net $\{\boldsymbol{h}_j\}_{j=1}^K$,
where $K \leq (2\operatorname{diam}(\mathcal{H})/\delta+1)^d = O(\epsilon^{-d/\alpha})$.

\textbf{Step 3 (Pointwise monotone approximation).}
For each centre $\boldsymbol{h}_j$, apply Lemma~\ref{lem:pwl} to
$g_j(\tau) := F_0(\tau|\boldsymbol{h}_j)$ on $[0,T^*]$ with tolerance $\epsilon/8$.
This yields a non-decreasing $\hat{g}_j \in [0,1]$ with $\hat{g}_j(0)=0$,
$O(\epsilon^{-1/\alpha})$ breakpoints, and
$\|\hat{g}_j - g_j\|_{\infty} < \epsilon/8$.

\textbf{Step 4 (Blending and error bound).}
Set $\hat{\mathcal{M}}(\tau,\boldsymbol{h}) := \sum_j w_j(\boldsymbol{h})\hat{g}_j(\tau)$.
Since each $\hat{g}_j$ is non-decreasing in $\tau$ and $w_j \geq 0$ with $\sum_j w_j = 1$,
$\hat{\mathcal{M}}$ is non-decreasing in $\tau$,
$\hat{\mathcal{M}} \in [0,1]$, and $\hat{\mathcal{M}}(0,\boldsymbol{h})=0$.

\emph{Core region $[0,T^*]\times\mathcal{H}$:}
\begin{align}
|\hat{\mathcal{M}}(\tau,\boldsymbol{h}) - F_0(\tau|\boldsymbol{h})|
&\leq \sum_j w_j\underbrace{|\hat{g}_j - g_j|}_{<\,\epsilon/8}
  + \sum_j w_j\underbrace{|F_0(\tau|\boldsymbol{h}_j) - F_0(\tau|\boldsymbol{h})|}_{\leq\, C\delta^\alpha(\|\boldsymbol{h}_j-\boldsymbol{h}\|/\delta)^\alpha}
\nonumber\\
&\leq \frac{\epsilon}{8} + C\delta^\alpha\underbrace{\sum_j w_j
  \bigl(\|\boldsymbol{h}_j-\boldsymbol{h}\|/\delta\bigr)^\alpha}_{\leq\,\Gamma}
\leq \frac{\epsilon}{8} + \frac{\epsilon}{8} = \frac{\epsilon}{4}. \label{eq:core}
\end{align}

\emph{Lower boundary} $\tau \in [0,\tau_{\min}]$:
By monotonicity and \eqref{eq:core}, $0 \leq \hat{\mathcal{M}}(\tau) \leq
\hat{\mathcal{M}}(\tau_{\min}) \leq F_0(\tau_{\min}) + \epsilon/4 < \epsilon/8 + \epsilon/4 = 3\epsilon/8$.
Since $F_0 \in [0, \epsilon/8]$, we have $|\hat{\mathcal{M}} - F_0| \leq 3\epsilon/8$.

\emph{Upper tail} $\tau > T^*$:
By monotonicity, $\hat{\mathcal{M}}(\tau) \geq \hat{\mathcal{M}}(T^*)
\geq F_0(T^*) - \epsilon/4 > 1 - 3\epsilon/8$.
Since $F_0 > 1 - \epsilon/8$ and $\hat{\mathcal{M}} \leq 1$,
$|\hat{\mathcal{M}} - F_0| \leq 3\epsilon/8$.

\smallskip
\emph{Stage A summary:}
$\|\hat{\mathcal{M}} - F_0\|_\infty < 3\epsilon/8$,
using $K \times O(\epsilon^{-1/\alpha}) = O(\epsilon^{-(1+d/\alpha)})$ parameters.

\medskip\noindent\textsc{Stage B: Lifting to the $\sigma$-MNN Architecture.}

\textbf{Step 5.}
We convert $\hat{\mathcal{M}}$ into $F^* = \sigma(\mathcal{M}_{\mathrm{net}})$
with $|F^* - \hat{\mathcal{M}}| < \epsilon/4 + \epsilon/8$.

\emph{(a) Perturbation.}
Define $\hat{\mathcal{M}}_\epsilon := (1-\epsilon/4)\hat{\mathcal{M}} + \epsilon/8 \in [\epsilon/8,\,1-\epsilon/8]$.
Since $\hat{\mathcal{M}} \in [0,1]$:
\begin{equation}
|\hat{\mathcal{M}}_\epsilon - \hat{\mathcal{M}}|
= |{-}\tfrac{\epsilon}{4}\hat{\mathcal{M}} + \tfrac{\epsilon}{8}|
\leq \tfrac{\epsilon}{8}.
\label{eq:perturb}
\end{equation}

\emph{(b) Scalar logit approximation.}
On $[a,b] := [\epsilon/8,\,1-\epsilon/8]$, the logit
$\operatorname{logit}(x) = \log(x/(1-x))$ is smooth with
$|\operatorname{logit}''(x)| \leq 64/\epsilon^2$.
By piecewise-linear interpolation on $M = O(\epsilon^{-3/2})$ equal subintervals
(Lemma~\ref{lem:pwl} applied to the logit on $[a,b]$):
$\|\hat{\ell} - \operatorname{logit}\|_{\infty,[a,b]} \leq 8/(\epsilon^2 M^2) \leq \epsilon$.
Since $\operatorname{Lip}(\sigma) = 1/4$:
\begin{equation}
|\sigma(\hat{\ell}(x)) - x|
= |\sigma(\hat{\ell}(x)) - \sigma(\operatorname{logit}(x))|
\leq \tfrac{1}{4}|\hat{\ell}(x) - \operatorname{logit}(x)|
\leq \tfrac{\epsilon}{4}.
\label{eq:logit_approx}
\end{equation}

\emph{(c) Composition.}
Define $\mathcal{M}_{\mathrm{net}}(\tau,\boldsymbol{h}) :=
\hat{\ell}(\hat{\mathcal{M}}_\epsilon(\tau,\boldsymbol{h}))$.
This is non-decreasing in $\tau$ and realisable as an MNN:
append to the Stage~A network an affine layer
($\hat{\mathcal{M}} \mapsto (1-\epsilon/4)\hat{\mathcal{M}} + \epsilon/8$; positive
$\tau$-weight since $1-\epsilon/4>0$) followed by a two-layer ReLU sub-network
with $O(\epsilon^{-3/2})$ non-negative weights implementing $\hat{\ell}$.
Set $F^* := \sigma(\mathcal{M}_{\mathrm{net}})$; it is non-decreasing and in $(0,1)$,
hence a valid CDF.

\emph{(d) Total error.}
By~\eqref{eq:logit_approx} with $x = \hat{\mathcal{M}}_\epsilon$:
\[
|F^* - F_0|
\;\leq\;
\underbrace{|\sigma(\mathcal{M}_{\mathrm{net}}) - \hat{\mathcal{M}}_\epsilon|}_{<\,\epsilon/4}
+ \underbrace{|\hat{\mathcal{M}}_\epsilon - \hat{\mathcal{M}}|}_{\leq\,\epsilon/8}
+ \underbrace{|\hat{\mathcal{M}} - F_0|}_{<\,3\epsilon/8}
\;<\; \frac{\epsilon}{4} + \frac{\epsilon}{8} + \frac{3\epsilon}{8}
= \frac{3\epsilon}{4} < \epsilon.
\]

\textbf{Parameter count.}
Stage A: $K\times O(\epsilon^{-1/\alpha}) = O(\epsilon^{-(1+d/\alpha)})$.
Stage B perturbation: $O(1)$.
Logit sub-network: $O(\epsilon^{-3/2})$.
For $\alpha \leq 1$ and $d \geq 1$, $1 + d/\alpha \geq 2 > 3/2$,
so Stage B is lower-order. Total: $N^* = O(\epsilon^{-(1+d/\alpha)})$.
\end{proof}

% ============================================================
\subsection{CDF--Intensity Equivalence}
% ============================================================

\begin{theorem}[CDF--Intensity Derivative Identity \verified{}]
\label{thm:cdf-intensity}
Let $\Lambda(\tau) := \int_0^\tau \lambda^*(s)\,ds$ and
$F^*(\tau) = 1 - e^{-\Lambda(\tau)}$.  Then:
\[
  \frac{dF^*(\tau)}{d\tau}
  = \lambda^*(\tau)\,(1 - F^*(\tau))
  = \lambda^*(\tau)\,e^{-\Lambda(\tau)}.
\]
In particular, the log-likelihood $\log p^*(\tau|\boldsymbol{h})
= \log(dF^*/d\tau)$ is finite whenever $\lambda^* > 0$.
\end{theorem}

\begin{proof}
Chain rule and FTC:
$\frac{d}{d\tau}[1 - e^{-\Lambda}] = e^{-\Lambda}\lambda^* = \lambda^*(1-F^*)$.
Log-likelihood finiteness follows since $\lambda^*>0$ and $e^{-\Lambda}>0$. \qed
\end{proof}

% ============================================================
\subsection{Numerical Superiority via Backward Error Analysis}
% ============================================================

We provide a rigorous \emph{backward error analysis} using condition numbers,
exposing a structural instability in intensity-based methods absent in CuFun.

\begin{definition}[Rounding Error Constant]
For $n$ elementary floating-point operations with unit roundoff $u < 1/n$:
$\gamma_n := nu/(1-nu)$.
\end{definition}

\begin{theorem}[Numerical Superiority via Condition Numbers]
\label{thm:numerical_superiority}
Let $u$ denote unit roundoff. Under Assumption~\ref{assumption:regularity},
let $\ell(\tau) := \log p^*(\tau)$ be the per-event log-likelihood.

\emph{(1) Intensity-based computation.}
For $N$-point trapezoidal quadrature, the computed log-likelihood
$\hat{\ell}_{\mathrm{int}}(\tau)$ satisfies the structural lower bound:
\[
|\hat{\ell}_{\mathrm{int}}(\tau) - \ell(\tau)|
\;\geq\; \frac{\tau u}{2} \min_i \lambda^*(s_i)  - O(N^{-2}).
\]
When $\log\lambda^*(\tau) \approx I(\tau) := \int_0^\tau\lambda^*(s)\,ds$
(near-stationary processes), the relative error diverges:
$|\hat{\ell}_{\mathrm{int}} - \ell|\,/\,|\ell| \to \infty$ as $\ell \to 0$.

\emph{(2) CuFun (CDF-based).}
For an MNN with $L$ layers and $W$ operations per layer,
the autodiff computed log-likelihood satisfies:
\[
|\hat{\ell}_{\mathrm{CDF}}(\tau) - \ell(\tau)| \;\leq\; \gamma_{2LW}\,|\ell(\tau)|.
\]
The condition number $\kappa_{\mathrm{CDF}} = O(1)$ is \emph{bounded independently of $N$}.
\end{theorem}

\begin{proof}
\textbf{Part (1).}
Using trapezoidal rule with $N$ evaluations at $s_i = i\tau/N$:
\[
\hat{I}_N(\tau) = \frac{\tau}{N}\sum_{i=0}^{N-1}
\frac{\lambda^*(s_i)(1+\delta_i) + \lambda^*(s_{i+1})(1+\delta_{i+1})}{2},
\quad |\delta_i| \leq u.
\]
The rounding contribution to $|\hat{I}_N - I|$ is bounded below by
$\frac{\tau u}{2}\min_i\lambda^*(s_i)$ up to the $O(N^{-2})$ truncation error.
The log-likelihood requires the subtraction
$\hat{\ell}_{\mathrm{int}} = \log\lambda^*(\tau) - \hat{I}_N(\tau)$.
When $a := \log\lambda^*(\tau) \approx b := I(\tau)$, this subtraction suffers
catastrophic cancellation: $|a - b|$ is small while each term has absolute error
$O(u|b|)$, so the \emph{relative} error in $a - b$ is $O(u|b|/|a-b|) \to \infty$
as $a \to b$.

\textbf{Part (2).}
CuFun computes $\ell = \log\sigma'(\mathcal{M}(\tau))\cdot\mathcal{M}'_\tau(\tau)$
via reverse-mode autodifferentiation. By Higham's standard backward error
analysis~\cite{higham2002accuracy} for $n = O(LW)$ elementary operations:
$|\hat{y} - y| \leq \gamma_n|y|$. No subtraction of nearly-equal quantities occurs:
$\log p^* = \log[\sigma'(\mathcal{M})\cdot\mathcal{M}'_\tau]$ combines positive terms,
and $\sigma'(z) = \sigma(z)(1-\sigma(z)) > 0$ everywhere by
Assumption~\ref{assumption:regularity}. The bound $\gamma_{2LW}$ is independent of
$N$ and of $\tau$.

\textbf{Comparison.}
For fixed $u, L, W$ and increasing $N$:
$|\hat{\ell}_{\mathrm{int}} - \ell| / |\hat{\ell}_{\mathrm{CDF}} - \ell|
\geq (\tau\min_i\lambda^*/2) / (\gamma_{2LW}|\ell|) \to \infty$ as $N\to\infty$,
since the numerator grows with cumulative summation error while the denominator
is fixed.
\end{proof}

% ============================================================
\subsection{Statistical Optimality: Minimax Rates}
% ============================================================

\begin{theorem}[Minimax Optimal Convergence]
\label{thm:consistency_rate}
Under Assumptions~\ref{assumption:foundation}--\ref{assumption:smoothness},
define the minimax risk:
$\mathcal{R}_n^* := \inf_{\hat{F}} \sup_{F_0 \in \mathcal{F}^\alpha}
\mathbb{E}\bigl[\|\hat{F}_n - F_0\|_{L^2}^2\bigr]$.

\emph{(Upper bound)} The CuFun MLE estimator achieves:
$\mathbb{E}\bigl[\|\hat{F}_n - F_0\|_{L^2}^2\bigr] = O\!\bigl(n^{-2\alpha/(2\alpha+d)}\bigr)$.

\emph{(Lower bound)} For any estimator based on $n$ i.i.d.\ samples:
$\sup_{F_0 \in \mathcal{F}^\alpha}\mathbb{E}\bigl[\|\hat{F}_n - F_0\|_{L^2}^2\bigr]
\geq C\cdot n^{-2\alpha/(2\alpha+d)}$.

Hence CuFun achieves the \emph{minimax optimal rate}.
\end{theorem}

\begin{proof}
\textbf{Upper bound.}
Decompose the MSE via bias--variance:
$\mathbb{E}[\|\hat{F}_n - F_0\|^2]
= \|F_m^* - F_0\|^2
+ \mathbb{E}[\|\hat{F}_n - F_m^*\|^2]$,
where $F_m^*$ is the best $m$-parameter approximator.

\emph{Approximation error.}
By Theorem~\ref{thm:universal_approx}, an $m$-parameter MNN achieves
$\|F_m^* - F_0\|_\infty \leq C_1 m^{-\alpha/d}$,
hence $\|F_m^* - F_0\|_{L^2}^2 \leq C_1^2 m^{-2\alpha/d}$.

\emph{Estimation error.}
Under Assumption~\ref{assumption:conditional_independence}, observations
$(\tau_i, \boldsymbol{h}_{i-1})$ are conditionally i.i.d.\ and the MLE minimizes
$\hat{R}_n(F) = -\frac{1}{n}\sum_i \log\frac{dF}{d\tau}(\tau_i|\boldsymbol{h}_{i-1})$.
By the Rademacher complexity bound for monotone neural networks
with $m$ parameters:
$\sup_{F \in \mathcal{F}_m} |\hat{R}_n(F) - R(F)| = O(\sqrt{m\log(nm)/n})$
with high probability.
Applying the standard localization argument, estimation error is $O(m\log(nm)/n)$.

\emph{Rate balancing.}
Set $m \sim (n/\log n)^{d/(2\alpha+d)}$; then
$m^{-2\alpha/d} \asymp m\log(nm)/n \asymp n^{-2\alpha/(2\alpha+d)}$
(up to log factors).

\textbf{Lower bound via Fano's inequality.}
We apply Fano's method~\cite{tsybakov2009introduction}.
Construct a $2\delta$-packing $\{F_v\}_{v \in \{0,1\}^K}$ in
$(\mathcal{F}^\alpha, \|\cdot\|_{L^2})$ with $K \geq \exp(c\delta^{-d/\alpha})$,
where $F_v = F_0 + \delta\sum_k v_k\psi_k$ and $\{\psi_k\}$ are disjoint
$\alpha$-H\"{o}lder bumps of unit $L^2$-norm.
By Assumptions~\ref{assumption:regularity}--\ref{assumption:smoothness}, the per-observation
KL divergence satisfies $\mathrm{KL}(P_{F_v}\|P_{F_{v'}}) \leq C\delta^2$.
Fano's inequality gives:
\[
\inf_{\hat{F}}\sup_{F_0}\mathbb{E}\|\hat{F}_n - F_0\|_{L^2}^2
\;\geq\; \frac{\delta^2}{4}\!\left(1 - \frac{nC\delta^2 + \log 2}{\log K}\right).
\]
Setting $\delta = n^{-\alpha/(2\alpha+d)}$ balances the bound, yielding
$\mathcal{R}_n^* \geq \frac{1}{8}\,n^{-2\alpha/(2\alpha+d)}$.
\end{proof}

\begin{remark}
Intensity-based MLE estimators introduce an additional bias from numerical
integration: $\mathbb{E}[(\hat{I}_N - I)^2] = O(\tau^4/N^2)$ per event.
This inflates the MSE to $O(n^{-2\alpha/(2\alpha+d)} + N^{-2})$, which matches
the minimax rate only when $N = \Omega(n^{\alpha/(2\alpha+d)})$, a non-trivial
computational overhead that grows with sample size.
\end{remark}
